\section{Hodge decomposition of the edge NTK}
\label{sec:theory}

\paragraph{Comparison graphs.} Items $1..n$ with features $x_i \in
\mathbb{R}^d$; a graph with oriented edges $(i,j)$, $i<j$; an edge flow $y
\in \mathbb{R}^E$, $y_{ij}$ the preference of $j$ over $i$ (a log-odds). With
the gradient $B:\mathbb{R}^V\to\mathbb{R}^E$, $(Bs)_{ij}=s_j-s_i$, and the
curl $C:\mathbb{R}^E\to\mathbb{R}^T$ over triangles, the Hodge decomposition
is $\mathbb{R}^E = \operatorname{im}B \oplus \operatorname{im}C^\top \oplus
\ker(BB^\top + C^\top C)$, with orthogonal projectors $P_{\text{grad}},
P_{\text{curl}}, P_{\text{harm}}$ of ranks $|V|-1$ (connected graph), the
number of independent triangles, and the first Betti number.

\paragraph{Three models and their edge kernels.} A preference model predicts
$g_{ij}$. The \emph{scalar} model $g_{ij}=f(x_j)-f(x_i)$; the
\emph{antisymmetric} model $g_{ij} = h([x_i,x_j]) - h([x_j,x_i])$; GPM $g_{ij}
= e(x_i)^\top R\, e(x_j)$ with $R$ block-skew. Each is an MLP in the NTK
parameterisation. Under gradient descent with step $\eta$ on
$\tfrac{1}{2E}\sum(g-y)^2$ the training residual evolves, in the linearised
regime, as $r_{t+1} = (I-\tfrac{\eta}{E}\Theta_E)\,r_t$, where $\Theta_E$ is
the NTK between edges. For the scalar model $\Theta_E = B\,\Theta_V B^\top$
with $\Theta_V$ the ordinary NTK between items; for the antisymmetric model,
with $z=[x_i,x_j]$ and $\tilde z=[x_j,x_i]$,
\begin{equation}
\Theta_E(e,e') = K(z,z') - K(z,\tilde z') - K(\tilde z,z') + K(\tilde z,\tilde z'),
\end{equation}
$K$ the closed-form ReLU NTK \cite{jacot2018neural} on concatenated inputs.
We define the \emph{Hodge rate} of a component as the mean eigenvalue of
$\Theta_E$ on its subspace, $\kappa_c = \operatorname{tr}(P_c\Theta_E P_c)
/ \dim_c$ --- the expected Rayleigh quotient of a random flow in that
component, hence its expected learning speed.

\paragraph{Proposition 1 (the scalar kernel is blind to cycles).} For any
cyclic flow $y_c \in \ker B^\top$, $\Theta_E\,y_c = B\Theta_V B^\top y_c = 0$.
The cyclic residual of a scalar model therefore never changes under gradient
descent on the squared loss, at any width and in any regime: it is not merely
unrepresentable but invisible to the gradient. Numerically $|\kappa_{\text{cyc}}|
\le \SpeScalarCycMax$ across all graphs.

\paragraph{Why the antisymmetric kernel prefers rankings.} Split the kernel on
concatenated inputs into a part additive in its two halves and an interaction
part. The additive part yields, after antisymmetrisation, a term of the form
$B\,\tilde K\,B^\top$ --- the scalar structure again --- whose action on cycles
vanishes by Proposition 1, while on gradient flows it is amplified by the
number of edges that share each item. Cycles are therefore learned only
through the interaction part. This predicts (i) $\kappa_{\text{grad}} \gg
\kappa_{\text{cyc}}$, (ii) a ratio that grows with the mean degree of the
comparison graph, and (iii) no reason for curl and harmonic cycles to differ.

\paragraph{Lemma 2 (i.i.d.\ noise is mostly cyclic).} If $z \sim
\mathcal{N}(0,\sigma^2 I_E)$, rotational invariance makes $\|P_{\text{grad}}
z\|^2/\|z\|^2 \sim \text{Beta}(\tfrac{|V|-1}{2}, \tfrac{|E|-|V|+1}{2})$, with
mean $(|V|-1)/|E|$. On a complete graph that is $2/n$: almost all noise lives
where a transitive-first learner looks last.
